%!TeX root=../main.tex
% Figura 2.4: parte do esquema do banco do kworkflow usada neste trabalho.
% Desenhada a partir de database/kwdb.sql (commit 2d8a04b). Mostra as tabelas
% do registro de execuções e a view statistics_report, que é a porta de
% entrada (gatilho INSTEAD OF INSERT) e de saída (kw report) desses dados.
% As tabelas do cronômetro de foco (timebox, tag) aparecem só pelo nome.
\begin{tikzpicture}[
    x=1cm, y=1cm,
    every node/.style={font=\footnotesize},
    tabela/.style={draw, rectangle split, rectangle split parts=2,
                   rectangle split part fill={black!12, white},
                   inner sep=3pt, font=\footnotesize},
    fora/.style={draw, dashed, text=black!55, draw=black!40,
                 inner sep=3pt, font=\footnotesize\ttfamily},
    visao/.style={draw, very thick, rounded corners=4pt, fill=black!4,
                  inner sep=4pt, align=center, font=\footnotesize},
    fk/.style={-{Latex[length=1.6mm]}, semithick},
    liga/.style={semithick, dashed},
  ]
  % Colunas: nome em \texttt à esquerda, tipo ou chave à direita
  \providecommand{\kwcolunas}{}\renewcommand{\kwcolunas}[1]{%
    \begin{tabular}{@{}l@{\hspace{1.2em}}r@{}}#1\end{tabular}}
  \providecommand{\kwcol}{}\renewcommand{\kwcol}[2]{\texttt{#1} & {\scriptsize #2}\\}

  \node[tabela] (event) at (0, 0) {%
    \texttt{event}
    \nodepart{two}
    \kwcolunas{\kwcol{id}{PK}\kwcol{date}{TEXT}\kwcol{time}{TEXT}}};

  \node[tabela] (executed) at (5.1, 0) {%
    \texttt{executed}
    \nodepart{two}
    \kwcolunas{\kwcol{id}{PK, FK}\kwcol{command\_label\_id}{FK}\kwcol{status\_id}{FK}%
             \kwcol{elapsed\_time\_in\_secs}{INTEGER}}};

  \node[tabela] (label) at (9.8, 0.8) {%
    \texttt{command\_label}
    \nodepart{two}
    \kwcolunas{\kwcol{id}{PK}\kwcol{name}{TEXT}}};

  \node[tabela] (status) at (9.8, -0.95) {%
    \texttt{status}
    \nodepart{two}
    \kwcolunas{\kwcol{id}{PK}\kwcol{name}{TEXT}}};

  \node[fora] (timebox) at (0, -2.1) {timebox, tag};

  % Chaves estrangeiras
  \draw[fk] (executed.west) -- node[above, font=\scriptsize] {mesmo \texttt{id}} (event.east);
  \draw[fk] ([yshift=1.5mm]executed.east) -- (label.west);
  \draw[fk] ([yshift=-2.5mm]executed.east) -- (status.west);
  \draw[fk, draw=black!40] (timebox.north) -- (event.south);

  % A view por cima
  \node[visao] (view) at (5.1, 2.6) {%
    view \texttt{statistics\_report}\\[1pt]
    {\ttfamily id, label\_name, status, date, time, elapsed\_time\_in\_secs}};
  \draw[liga] ([xshift=4mm]view.south west) -- (event.north);
  \draw[liga] (view.south) -- (executed.north);
  \draw[liga] ([xshift=-4mm]view.south east) -- (label.north);
\end{tikzpicture}
